% ATTEMPT_STATUS: unresolved
% TOP_PROBLEM: {"problemId": "problem.long-range-order-in-the-quantum-heisenberg-ferromagnet", "canonicalTitle": "Long-Range Order in the Quantum Heisenberg Ferromagnet", "releaseRank": 202, "primaryDomain": "mathematical_physics", "primaryDomainLabel": "Mathematical physics", "displayStatus": "open_with_solved_subcases", "releaseStatus": "open_with_solved_subcases"}
% !TeX program = lualatex
% Complete problem section copied verbatim from research_batch_201_202.tex.
\documentclass[11pt]{article}
\usepackage[margin=25mm]{geometry}
\usepackage{fontspec}
\setmainfont{Latin Modern Roman}
\usepackage{amsmath,amssymb,mathtools}
\usepackage{unicode-math}
\setmathfont{Latin Modern Math}
\usepackage{enumitem,needspace}
\usepackage{xurl,hyperref}
\hypersetup{colorlinks=true,urlcolor=blue}
\setcounter{secnumdepth}{1}
\setlength{\parindent}{0pt}
\setlength{\parskip}{5pt}
\setlength{\emergencystretch}{3em}
\setlist[itemize]{leftmargin=3em,itemsep=5pt,topsep=4pt,parsep=0pt}
\allowdisplaybreaks
\newcommand{\N}{\mathbb{N}}
\newcommand{\Z}{\mathbb{Z}}
\newcommand{\Q}{\mathbb{Q}}
\newcommand{\R}{\mathbb{R}}
\newcommand{\C}{\mathbb{C}}
\newcommand{\F}{\mathbb{F}}
\newcommand{\A}{\mathbb{A}}
\newcommand{\PP}{\mathbb P}
\newcommand{\E}{\mathbb{E}}
\newcommand{\eps}{\varepsilon}
\newcommand{\dd}{\,\mathrm d}
\newcommand{\sref}[2]{\hyperlink{source:#1:#2}{[S#2]}}
\newcommand{\eref}[2]{\hyperlink{extra:#1:#2}{[E#2]}}
% Turn the next switch off for a shorter reading copy. The quotations preserve
% every exactTarget field, including qualifications omitted from a short summary.
\newif\ifshowcatalogue
\showcataloguetrue
\newcommand{\cataloguescope}[1]{\ifshowcatalogue
	\paragraph{Exact scope in the supplied catalogue.}
	\begin{quote}\small #1\end{quote}\fi}
\begin{document}
\setcounter{section}{201}
	% ================================================================
	% problemId: problem.long-range-order-in-the-quantum-heisenberg-ferromagnet
	% source releaseRank: 202; source releaseStatus: open_with_solved_subcases
	% exactTarget SHA-256: 6b6155a73c536e1df98675d36c115216d36a32dc833614a37b7c949626ad07c8
	\clearpage
	\section{\texorpdfstring{Long-Range Order in the Quantum Heisenberg Ferromagnet}{Long-Range Order in the Quantum Heisenberg Ferromagnet}}\label{202}
	\subsection{Definitions and mathematical statement}
	At each site of $\Z^d$, spin-$S$ operators act on $\C^{2S+1}$, satisfy the angular-momentum commutation relations, and have $\mathbf S_x^2=S(S+1)I$. The finite-volume Hamiltonian is
	$H_\Lambda=-\sum_{\langle x,y\rangle\subset\Lambda}\mathbf S_x\cdot\mathbf S_y$.
	An infinite-volume equilibrium state is understood in the quantum Gibbs/KMS sense of the cited model. For each $d\ge3$ and positive half-integer $S$, the assertion is that some finite $\beta_0(d,S)$ satisfies
	\[
	\forall\beta\ge\beta_0\text{ finite},\quad
	\exists\text{ translation-invariant zero-field equilibrium state }\omega_\beta,
	\quad\omega_\beta(\mathbf S_0)\ne0.
	\]
	The state is allowed to break rotational symmetry; its rotation-invariant mixture need not have nonzero magnetization. The coupling is fixed to one.
	\par\smallskip\textit{Formulation and concept references:} \sref{202}{1}.
	\cataloguescope{For the nearest-neighbor isotropic spin-S quantum Heisenberg ferromagnet on Z\textasciicircum{}d with Hamiltonian H=-sum\_\{<x,y>\} S\_x dot S\_y, assert: for every integer d>=3 and every S in\{1/2,1,3/2,...\}, there is a finite beta\_0(d,S)>0 such that for each finite beta>=beta\_0 the infinite-volume zero-field model admits a translation-invariant equilibrium Gibbs state with nonzero spin expectation. Nonzero expectation refers to spontaneous magnetization in a symmetry-breaking state, not the rotation-invariant mixture. Coupling J=1 fixes temperature units; no pointwise-correlation equivalence is asserted.}
	\subsection{Short English statement}
	At sufficiently low positive temperature, must the quantum Heisenberg ferromagnet in every dimension at least three possess an equilibrium state with spontaneous magnetization?
	% BEGIN RESEARCH ADDITION 202
	\subsection{Research attempt}
	
	\paragraph{Current frontier and scope (24 September 2026).}
	The target is a symmetry-breaking equilibrium state at each sufficiently
	low \emph{positive} temperature, not a magnetized finite-volume zero-field
	trace and not merely an ordered ground state. The original source
	\sref{202}{1} distinguishes the quantum ferromagnet from the classical
	and quantum-antiferromagnetic problems. For the present isotropic model,
	Correggi--Giuliani--Seiringer prove the leading three-dimensional
	low-temperature free energy for every fixed spin. With the ground
	energy shifted to zero, their Theorem 2.1 gives
	\[
	f^0(\beta,S)\sim
	-\frac{\zeta(5/2)}{8\pi^{3/2}}S^{-3/2}\beta^{-5/2}.
	\]
	This is a published free-energy theorem, not a theorem about spontaneous
	magnetization.\ \eref{202}{1}
	
	Two method cautions matter. Speer proves failure of the standard
	reflection-positivity property in specified finite-volume spin-$1/2$
	regimes, so that property cannot be imported wholesale from other spin
	models; this does not rule out every possible infrared argument.
	\eref{202}{2} The January 2025 revision of Heson--Starr--Thornton studies
	energy-ordering reversals on even spin rings, with a largely numerical
	part and a rigorous singlet-sector result. It does not establish the
	all-dimensional lattice target.\ \eref{202}{5}
	None of the sources inspected supplies the full assertion in this entry.
	The calculations below are independent deductions; any claim of their
	novelty would require a separate literature comparison.
	
	\paragraph{Attack map.}
	The first route is to differentiate a spin-wave free-energy estimate in
	a field. It needs error control on the scale of the field itself, not
	only an asymptotic at zero field. The second route resolves the
	finite-volume trace into total-spin sectors, removes the rotational
	zero mode, and estimates the loss of total spin. It needs a
	volume-uniform bound on that loss. The third route, available directly
	for $S=1/2$, uses weighted interchange cycles and needs positive density
	of macroscopic cycles in that \emph{weighted} model. These routes are
	related, but none of the missing estimates is a theorem supplied by a
	formal spin-wave substitution.
	
	The second route is pursued first. It produces a precise state-selection
	lemma and a conditional reduction. I then attempt a strong free-boson
	comparison that would close the reduction, and disprove that proposed
	comparison. A spatial twist supplies an obstruction even on the tori
	approximating the requested lattice, not just on an unrelated graph.
	
	\paragraph{Attempt A: separate total-spin loss from global rotation.}
	Let $\Lambda_L=(\Z/L\Z)^d$, $L\geq3$, with its nearest-neighbor torus
	graph, and write $N=L^d$. Each unordered bond is counted once. Put
	\begin{equation}\label{ra202:operators}
		H^0_L=H_L+dNS^2I
		=\sum_{\langle x,y\rangle}(S^2I-\mathbf S_x\cdot\mathbf S_y),
		\quad \mathbf J=\sum_x\mathbf S_x,\quad M=J^3,
		\quad \mathcal J=\frac{\sqrt{I+4\mathbf J^2}-I}{2}.
	\end{equation}
	The largest eigenvalue of $\mathbf S_x\cdot\mathbf S_y$ is $S^2$,
	by addition of two spins, so $H^0_L\geq0$. The operator $\mathcal J$
	has eigenvalue $j$ on a total-spin-$j$ representation, with
	$0\leq j\leq SN$. Define
	\[
	D=SN I-\mathcal J.
	\]
	All three operators $H^0_L$, $\mathcal J$, and $M$ commute. The
	quantity $D$ measures loss of \emph{spin length}, whereas $SN-M$ also
	counts a uniform rotation away from the chosen third axis. This
	difference cannot be ignored: rotational invariance gives
	\begin{equation}\label{ra202:zeromode}
		\langle M\rangle_{\beta,0,L}=0,
		\qquad \langle SN-M\rangle_{\beta,0,L}=SN
	\end{equation}
	for every finite volume and every $\beta$, including arbitrarily large
	$\beta$. A claim that the latter occupation is small in the zero-field
	trace is therefore false before any limiting argument is attempted.
	
	\textbf{Lemma 202.1 (finite-volume field selection).}
	For any finite $\mathrm{SU}(2)$-invariant Hamiltonian on $N$ spin-$S$
	sites, and for $\beta,h>0$, define the thermal state using
	$H^0-hM$. Then
	\begin{equation}\label{ra202:selection}
		\frac{\langle M\rangle_{\beta,h,L}}{N}
		\geq \frac{\langle\mathcal J\rangle_{\beta,0,L}}{N}
		-\frac{1}{N(e^{\beta h}-1)}.
	\end{equation}
	No reflection positivity or sign assumption on that Hamiltonian is
	needed for this finite-dimensional statement.
	
	\emph{Proof.}
	Decompose the Hilbert space as
	$\bigoplus_j(V_j\otimes\mathcal M_j)$, where $V_j$ is the spin-$j$
	irreducible representation. Schur's lemma gives
	$H^0=\bigoplus_j(I_{V_j}\otimes H_j)$. Set
	$w_j=\operatorname{Tr}_{\mathcal M_j}e^{-\beta H_j}\geq0$ and $x=\beta h$.
	The spin character is
	\[
	\chi_j(x)=\sum_{m=-j}^j e^{xm}
	=e^{jx}\frac{1-e^{-(2j+1)x}}{1-e^{-x}},
	\qquad
	\frac{\chi_j'(x)}{\chi_j(x)}
	=j-\frac1{e^x-1}+\frac{2j+1}{e^{(2j+1)x}-1}
	\geq j-\frac1{e^x-1}.
	\]
	The zero-field law of $j$ is proportional to $(2j+1)w_j$. The field
	reweights it by $a_j(x)=\chi_j(x)/(2j+1)$. For fixed $x>0$ this is
	nondecreasing in $j$ along the common integer or half-integer parity:
	passing from $j$ to $j+1$ adds two endpoints whose average is
	$\cosh((j+1)x)$, no smaller than the previous average. For any probability
	weights $p_j$,
	\[
	\operatorname{Cov}_p(j,a_j)
	=\frac12\sum_{j,k}p_jp_k(j-k)(a_j-a_k)\geq0.
	\]
	Consequently the mean of $j$ cannot decrease under this reweighting.
	Averaging the character derivative gives
	$\langle M\rangle_{\beta,h,L}\geq
	\langle\mathcal J\rangle_{\beta,0,L}-(e^{\beta h}-1)^{-1}$.
	Division by $N$ proves the assertion.\hfill$\square$
	
	\paragraph{The conditional reduction, including the equilibrium-state step.}
	For fixed $d,S,\beta$, consider the quantitative condition
	\begin{equation}\label{ra202:HW}
		\liminf_{L\to\infty}
		\frac{\langle\mathcal J\rangle_{\beta,0,L}}{N}
		\geq m_\beta>0. \tag{HW}
	\end{equation}
	\textbf{Proposition 202.2.} Condition \eqref{ra202:HW} implies the
	existence of a translation-invariant zero-field $\beta$-KMS state with
	$\omega_\beta(S_0^3)\geq m_\beta$.
	
	\emph{Proof.}
	Choose $h_L=N^{-1/2}$. Equation \eqref{ra202:selection}, together with
	$e^x-1\geq x$, gives
	\[
	\langle S_0^3\rangle_{\beta,h_L,L}
	=\frac{\langle M\rangle_{\beta,h_L,L}}{N}
	\geq\frac{\langle\mathcal J\rangle_{\beta,0,L}}{N}
	-\frac1{\beta\sqrt N}.
	\]
	The equality uses torus translation invariance. By successive extraction
	on local matrix algebras, these states have a subsequence converging on
	every local observable. Positivity and consistency give a state on the
	quasi-local algebra. It is translation invariant and has the asserted
	spin expectation.
	
	Here is why the limit is an \emph{equilibrium} state, rather than merely
	a positive functional. For bounded finite-range interactions the
	finite-volume dynamics converge in norm on local observables, uniformly
	on compact real-time intervals; this is the locality estimate proved in
	\eref{202}{4}, Section 3.1, Theorem 3.1. Periodic boundary bonds make no
	difference on any fixed local algebra in this limit: their distance
	from that algebra tends to infinity in the same locality estimate.
	Since $[H_L^0,M]=0$, adding $-h_LM$ composes the zero-field dynamics
	with an on-site spin rotation of angle $-h_Lt$. On an observable
	supported on $X$, its difference from the identity is bounded by
	$2S|X|\,|h_Lt|\|A\|$. Thus these field-dependent dynamics also converge
	to the \emph{zero-field} dynamics, denoted $\alpha_t$.
	
	To justify passage of the KMS condition without asserting convergence
	at imaginary time for an unsmeared local observable, fix $a>0$ and set
	\[
	B_{L,a}=\sqrt{a/\pi}\int_{\R} e^{-at^2}\alpha_t^{L,h_L}(B)\dd t.
	\]
	These elements and their analytic translates converge in norm to the
	corresponding $B_a$ and $\alpha_z(B_a)$, because
	\[
	\alpha_z^{L,h_L}(B_{L,a})
	=\sqrt{a/\pi}\int_{\R} e^{-a(t-z)^2}\alpha_t^{L,h_L}(B)\dd t
	\]
	has an integrable, volume-independent majorant for each fixed $z$.
	The finite-volume trace identity
	\[
	\omega_L\bigl(A\alpha_{i\beta}^{L,h_L}(B_{L,a})\bigr)
	=\omega_L(B_{L,a}A)
	\]
	therefore passes to the limit. Gaussian-smoothed observables form the
	usual dense analytic test set for the KMS condition; allowing local
	$A,B$ and then their norm limits proves that the limiting state is
	$\beta$-KMS. This establishes the required equilibrium property at the
	fixed finite $\beta$.\hfill$\square$
	
	In particular, it would suffice to prove, for each fixed $d\geq3$ and
	$S$, some $\rho_{d,S}(\beta)<S$ such that
	\begin{equation}\label{ra202:defectgoal}
		\limsup_{L\to\infty}\frac{\langle D\rangle_{\beta,0,L}}{N}
		\leq\rho_{d,S}(\beta)
		\qquad\text{for every finite }\beta\geq\beta_0(d,S).
	\end{equation}
	Then \eqref{ra202:HW} holds with $m_\beta=S-\rho_{d,S}(\beta)$.
	This is a \emph{sufficient condition}, not an unproved equivalence
	between a pointwise correlation limit and the original statement.
	Notice the order and the quantifiers: $\beta$ is fixed during the
	thermodynamic limit; a ground-state limit taken first does not suffice,
	and no monotonicity in $\beta$ has been assumed.
	
	\paragraph{Attempt B: spin waves suggest a bound, but do not prove it.}
	For the quadratic spin-wave model with the coupling normalization of
	this entry, the nonzero momentum energies are
	\[
	\epsilon_S(k)=2S\sum_{i=1}^d(1-\cos k_i),\qquad
	\rho_{\rm sw}(\beta)=
	\int_{[-\pi,\pi]^d}\frac{1}{e^{\beta\epsilon_S(k)}-1}
	\frac{\dd k}{(2\pi)^d}.
	\]
	The elementary bound $\epsilon_S(k)\geq cS|k|^2$ gives, for $d>2$ and
	$\beta S\geq1$,
	\begin{equation}\label{ra202:bosedensity}
		\rho_{\rm sw}(\beta)
		\leq C_d(\beta S)^{-d/2}
		\int_{\R^d}\frac{\dd q}{e^{c|q|^2}-1}
		\leq C_d'(\beta S)^{-d/2}.
	\end{equation}
	The integral is finite precisely because near zero its radial
	integrand behaves like $r^{d-3}$. This calculation explains the
	favorable dimension threshold of this route. It estimates only the
	free Bose integral, not the interacting spin expectation.
	
	A strong enough proposed estimate would be
	\begin{equation}\label{ra202:strongdefect}
		\limsup_{L\to\infty}\frac{\langle D\rangle_{\beta,0,L}}{N}
		\leq C_{d,S}(\beta S)^{-d/2}
		\quad(\beta S\geq1). \tag{SD}
	\end{equation}
	If \eqref{ra202:strongdefect} were proved, one could take
	\[
	\beta_0(d,S)=S^{-1}\max\left\{1,
	(2C_{d,S}/S)^{2/d}\right\}
	\]
	and obtain a state with $\omega_\beta(S_0^3)\geq S/2$ for every finite
	$\beta\geq\beta_0$. Equation \eqref{ra202:strongdefect} is not proved
	here. The exact Holstein--Primakoff model has an occupation constraint
	and attractive interaction terms; the free-energy analysis in
	\eref{202}{1}, Section 3, does not discard them as an operator
	inequality giving \eqref{ra202:strongdefect}.
	
	I tried a sharper comparison that would have made differentiation
	legitimate. For a finite torus, omit $k=0$ and let
	$\epsilon_1=\min_{k\ne0}\epsilon_S(k)$. The candidate was
	\begin{equation}\label{ra202:falsemgf}
		\left\langle z^D\right\rangle_{\beta,0,L}
		\ \stackrel{?}{\leq}\
		\prod_{k\ne0}\frac{1-e^{-\beta\epsilon_S(k)}}
		{1-ze^{-\beta\epsilon_S(k)}},
		\qquad 1\leq z<e^{\beta\epsilon_1}.
	\end{equation}
	Both sides equal one at $z=1$. A valid bound of this form in a
	right-neighborhood of one would bound the first derivative, hence
	$\langle D\rangle$, by the finite-volume free-boson number. The latter
	per site tends to \eqref{ra202:bosedensity} for $d>2$: near zero the
	Riemann sums are controlled by the summable singularity $|k|^{-2}$.
	However, the \emph{full interval} and finite-volume uniformity asserted
	in \eqref{ra202:falsemgf} are false, as shown next. Differentiating a
	false stronger statement cannot establish the desired weaker one.
	
	\paragraph{Stress test I: a slow twist defeats energy coercivity.}
	\textbf{Lemma 202.3 (low total spin at vanishing energy density).}
	For each fixed $d,S$ and all sufficiently large $L$, there is a joint
	eigenvector of $H_L^0$ and $\mathcal J$ with
	\begin{equation}\label{ra202:lowspin}
		j\leq N^{3/4},\qquad
		E\leq\frac{NS^2(1-\cos(2\pi/L))}{1-S/\sqrt N}.
	\end{equation}
	
	\emph{Proof.}
	For every site set
	\[
	\mathbf n_x=(\cos(2\pi x_1/L),\sin(2\pi x_1/L),0),
	\qquad \Psi_L=\bigotimes_x |S,\mathbf n_x\rangle,
	\]
	where $|S,\mathbf n\rangle$ is the spin coherent state with expectation
	$S\mathbf n$. These states exist by rotating a highest-weight vector.
	Their product expectations factor at different sites, and
	$\sum_x\mathbf n_x=0$. Only bonds in the first coordinate direction
	cost energy. Direct calculation therefore gives
	\begin{align}
		\langle H_L^0\rangle_{\Psi_L}
		&=NS^2(1-\cos(2\pi/L))\leq 2\pi^2S^2N/L^2,
		\label{ra202:twistenergy}\\
		\langle\mathbf J^2\rangle_{\Psi_L}
		&=NS(S+1)+S^2\sum_{x\ne y}\mathbf n_x\cdot\mathbf n_y=NS.
		\label{ra202:twistspin}
	\end{align}
	Since $\mathcal J^2\leq\mathbf J^2$, Cauchy--Schwarz yields
	$\langle\mathcal J\rangle_{\Psi_L}\leq\sqrt{NS}$ and hence
	$\langle D\rangle_{\Psi_L}\geq SN-\sqrt{NS}$.
	Already this disproves any operator inequality $H_L^0\geq cD$ with
	$c>0$ independent of $L$.
	
	To obtain an eigenvector rather than only a trial state, let
	$P=\mathbf1_{\{\mathcal J\leq N^{3/4}\}}$. Markov's inequality and
	\eqref{ra202:twistspin} imply
	\[
	\|P\Psi_L\|^2\geq1-S/\sqrt N.
	\]
	The operator $P$ commutes with $H_L^0$, and $H_L^0\geq0$; therefore
	\[
	\frac{\langle P\Psi_L,H_L^0P\Psi_L\rangle}{\|P\Psi_L\|^2}
	\leq\frac{\langle\Psi_L,H_L^0\Psi_L\rangle}{1-S/\sqrt N}.
	\]
	The minimum eigenvalue in the range of $P$ is at most this Rayleigh
	quotient. Simultaneous diagonalization inside that range proves
	\eqref{ra202:lowspin}.\hfill$\square$
	
	The lemma shows why a total energy-density bound cannot on its own
	force extensive total spin: states with $j/N\to0$ can have $E/N\to0$.
	It does not show that these states dominate a fixed-temperature trace.
	Their entropy and their relative thermal weights are exactly what an
	argument for \eqref{ra202:defectgoal} must still control. Nor does it
	exclude a bound with a factor $L^{-2}$, of the type used on boxes in
	\eref{202}{1}, Proposition 5.2; that vanishing factor is important.
	
	\paragraph{Stress test II: the proposed fugacity comparison fails on the target tori.}
	On these tori the smallest nonzero quadratic spin-wave energy is
	\[
	\epsilon_1=2S(1-\cos(2\pi/L)).
	\]
	For the joint eigenvector in \eqref{ra202:lowspin}, let $D_*=SN-j$.
	For fixed $S$ and sufficiently large $N$,
	\begin{equation}\label{ra202:halfgap}
		\frac{E}{\epsilon_1D_*}
		\leq\frac{1}
		{2(1-S/\sqrt N)(1-N^{-1/4}/S)}<\frac34.
	\end{equation}
	Now fix one such torus and put $z=e^{3\beta\epsilon_1/4}$.
	It lies in the purported domain of \eqref{ra202:falsemgf}. Since
	$H_L^0\geq0$, its partition function is at most $(2S+1)^N$, while the
	single eigenvector contributes
	\[
	\left\langle z^D\right\rangle_{\beta,0,L}
	\geq (2S+1)^{-N}
	\exp\{\beta(3\epsilon_1D_*/4-E)\}
	\longrightarrow\infty
	\quad\text{as }\beta\to\infty.
	\]
	In contrast, every factor on the right of \eqref{ra202:falsemgf} tends
	to one, since $\epsilon_S(k)-3\epsilon_1/4\geq\epsilon_1/4$ and
	there are only finitely many factors. This is an exact contradiction
	to that finite-volume comparison, for every fixed $d\geq3$ as well as
	for other dimensions. It is \emph{not} a contradiction to
	\eqref{ra202:strongdefect}: the test first fixes the volume and then
	lets $\beta$ grow, and its fugacity moves far away from one.
	A much weaker comparison of first derivatives after the correct
	thermodynamic limit is still conceivable.
	
	An exact four-spin check makes the mechanism transparent. A numerical
	scan of magnetization blocks on cycles of $4,6,8,10$ sites detected
	states with $E<\epsilon_1D$; the following calculation
	removes any reliance on that experiment. For a spin-$1/2$ four-cycle,
	set $\mathbf A=\mathbf S_1+\mathbf S_3$ and
	$\mathbf B=\mathbf S_2+\mathbf S_4$. Then
	\[
	H^0=1-\mathbf A\cdot\mathbf B,
	\qquad E=1-\tfrac12\{j(j+1)-a(a+1)-b(b+1)\}.
	\]
	Taking opposite pairs to be singlets gives $a=b=j=0$, so $E=1$ and
	$D=2$, whereas the nonzero quadratic energies are $1,1,2$. The complete
	spin-addition calculation gives
	\[
	Z=5+7e^{-\beta}+3e^{-2\beta}+e^{-3\beta},\qquad
	\langle z^D\rangle=
	\frac{5+(6z+z^2)e^{-\beta}+3ze^{-2\beta}+z^2e^{-3\beta}}{Z}.
	\]
	For $z=e^{3\beta/4}$ the same contradiction is explicit. This small
	cycle is a diagnostic for the comparison, not a counterexample to
	magnetization on $\Z^d$; the preceding twist argument is what carries
	the comparison failure to the relevant large tori.
	
	\paragraph{Attempt C: a separate loop criterion, with the correct spin restriction.}
	For $S=1/2$, let $T_{xy}$ interchange the two tensor factors. Then
	\[
	T_{xy}=2\mathbf S_x\cdot\mathbf S_y+\tfrac12I,
	\qquad H_L^0=\tfrac12\sum_{\langle x,y\rangle}(I-T_{xy}).
	\]
	Place independent Poisson transpositions of rate $1/2$ on each edge
	during $[0,\beta]$, and let $\pi$ be their chronological product.
	Expansion of the exponential gives
	$e^{-\beta H_L^0}=\mathbb E_{\rm Pois}T_\pi$.
	In a spin basis, $\operatorname{Tr}T_\pi=2^{c(\pi)}$, since a basis
	configuration fixed by $\pi$ is constant on each permutation cycle.
	Thus the probability law relevant to this quantum model is
	\begin{equation}\label{ra202:looplaw}
		\dd\mathbb P_2(\pi)=
		\frac{2^{c(\pi)}}{\mathbb E_{\rm Pois}2^{c(\pi)}}
		\dd\mathbb P_{\rm Pois}(\pi),
	\end{equation}
	not the unweighted interchange law. This is the T\'oth-type
	representation discussed in \eref{202}{3}, Theorem 3.2 and Section 7.1;
	the rate $1/2$ here accounts for the coupling convention of this entry.
	
	Conditional on the cycle lengths $\ell_i$, independent signs
	$\sigma_i\in\{-1/2,1/2\}$ give
	$M=\sum_i\ell_i\sigma_i$ in the diagonal trace. Averaging the signs
	therefore proves the exact identity
	\begin{equation}\label{ra202:loopsquare}
		\langle M^2\rangle_{\beta,0,L}
		=\frac14\mathbb E_2\sum_i\ell_i^2.
	\end{equation}
	For example, if some $\eps,\delta>0$ satisfy
	\begin{equation}\label{ra202:loopgoal}
		\liminf_{L\to\infty}\mathbb E_2
		\left[\frac1N\sum_{i:\ell_i\geq\eps N}\ell_i\right]\geq\delta,
	\end{equation}
	then \eqref{ra202:loopsquare} gives
	$\liminf N^{-2}\langle M^2\rangle\geq\eps\delta/4$.
	On a spin-$j$ sector, $M^2\leq j^2\leq(SN+1)j$, so
	\[
	\frac{\langle\mathcal J\rangle}{N}
	\geq\frac{\langle M^2\rangle}{N(SN+1)}.
	\]
	For $S=1/2$, condition \eqref{ra202:HW} follows with
	$m_\beta=\eps\delta/2$, and Proposition 202.2 supplies the required
	symmetry-breaking state at that $\beta$. Thus the bridge from the
	cycle criterion to the \emph{state formulation} has been proved, not
	assumed as a pointwise-correlation equivalence.
	
	The attempted percolation shortcut fails. An added transposition can
	join two cycles or split one cycle; two successive transpositions on
	the same two vertices can undo one another. The event of a long cycle
	is not an increasing event of the Poisson bridges. Moreover the weight
	$2^{c(\pi)}$ changes in opposite directions under joining and splitting.
	A theorem about ordinary percolation or unweighted interchange therefore
	does not automatically prove \eqref{ra202:loopgoal}. The missing weighted
	cycle estimate is not supplied here.
	
	The larger-spin target has an additional obstruction to this shortcut.
	For $S>1/2$, the bilinear interaction is not an affine function of the
	full-site transposition. For example, at $S=1$ the transposition has
	the same eigenvalue on pair-spin $0$ and $2$, while
	$\mathbf S_x\cdot\mathbf S_y$ has the different eigenvalues $-2$ and
	$1$ there. Simply changing $2^{c(\pi)}$ to $(2S+1)^{c(\pi)}$ would
	change the Hamiltonian. The all-$S$ state-selection reduction remains
	valid, but the displayed cycle criterion directly addresses only
	$S=1/2$.\ \eref{202}{3}, Sections 3.1 and 7.
	
	\paragraph{Final stress test: why a free-energy asymptotic does not close the gap.}
	If $f_\beta(h)$ denotes free energy per site, magnetization is related
	to its one-sided slope at zero, not just its value there. Uniform
	small errors in values need not control this slope. The elementary
	concave, even functions
	\[
	q_\delta(h)=-S\bigl(\sqrt{h^2+\delta^2}-\delta\bigr)
	\]
	satisfy $q_\delta(0)=q_\delta'(0)=0$ and approximate $-S|h|$ uniformly
	within $S\delta$. Nevertheless the limiting right derivative is
	$-S$, not zero. Thus even a very accurate value approximation may
	lose the relevant one-sided derivative. In particular, the zero-field
	asymptotic in \eref{202}{1} cannot be differentiated into spontaneous
	magnetization without a separate error estimate at field scale.
	
	There is also a mesoscopic obstruction. Nearly aligned spins within
	every small box need not align their directions between boxes: the
	slow twist above has this property while its total spin is small.
	A boxwise spin-wave estimate must be supplemented by a bound on the
	relative orientations and their entropy. The selection lemma does
	not supply that bound. Reflection positivity, monotonicity of cycle
	lengths, or a bosonic partition-function domination is not assumed
	at the unresolved step.
	
	\paragraph{Outcome.}
	\textbf{Outcome: Conditional reduction.}
	
	The finite-volume field-selection inequality and its zero-field KMS
	limit prove that \eqref{ra202:defectgoal}, for every sufficiently large
	fixed finite $\beta$, would imply the exact state-existence target for
	each $d,S$. A macroscopic-cycle condition gives a second sufficient
	criterion for $S=1/2$. These are proved implications with explicit
	normalizations and limit orders; the required defect or cycle bound
	is still unproved.
	
	The counterexample search has a concrete, narrower outcome: coherent
	twists and an exact four-spin calculation disprove the proposed full
	fugacity comparison with zero-mode-removed free bosons. The twist also
	rules out volume-independent energy coercivity in total-spin loss.
	Neither construction is a counterexample to the original conjecture.
	What remains is a genuinely interacting, thermodynamic estimate of
	total-spin loss or weighted macroscopic cycles, not a change of the
	original all-dimensional, all-spin, positive-temperature statement.
	% END RESEARCH ADDITION 202
	\subsection{Sources}
	\begin{itemize}\raggedright
		\item[\textnormal{[S1]}]\hypertarget{source:202:1}{} Long range order for the quantum Heisenberg model, contributed27January1999.
		\url{https://web.math.princeton.edu/~aizenman/OpenProblems_MathPhys/9901.HeisenbergFerr.html}.
		{\small\textit{Catalogue formulation source.}}
		% BEGIN RESEARCH REFERENCES 202
		\item[\textnormal{[E1]}]\hypertarget{extra:202:1}{}
		M. Correggi, A. Giuliani, and R. Seiringer, \emph{Validity of the
			spin-wave approximation for the free energy of the Heisenberg
			ferromagnet}, Communications in Mathematical Physics 339 (2015),
		279--307, Theorem 2.1, Section 3, and Proposition 5.2.
		\url{https://arxiv.org/pdf/1312.7873v2};
		\url{https://doi.org/10.1007/s00220-015-2402-0}.
		{\small\textit{Published result; arXiv version 2 is dated 10 August 2019.
				Consulted 24 September 2026.}}
		\item[\textnormal{[E2]}]\hypertarget{extra:202:2}{}
		E. R. Speer, \emph{Failure of reflection positivity in the quantum
			Heisenberg ferromagnet}, Letters in Mathematical Physics 10 (1985),
		41--47, abstract, Introduction, and the finite-volume counterexamples.
		\url{https://sites.math.rutgers.edu/~speer/homepage/pubs/LMP_10_41.pdf}.
		{\small\textit{Its specified regimes are not replaced here by a claim
				about every temperature or every possible reflection. Consulted
				24 September 2026.}}
		\item[\textnormal{[E3]}]\hypertarget{extra:202:3}{}
		D. Ueltschi, \emph{Random loop representations for quantum spin systems},
		Journal of Mathematical Physics 54 (2013), 083301; author manuscript
		arXiv:1301.0811, Theorem 3.2, the spin/loop correspondence of Section 3,
		and Section 7.1, especially equation (7.1).
		\url{https://arxiv.org/pdf/1301.0811}.
		{\small\textit{The spin-$1/2$ ferromagnetic representation is used, not
				the different models for which this paper proves long-range order.
				Consulted 24 September 2026.}}
		\needspace{9\baselineskip}
		\item[\textnormal{[E4]}]\hypertarget{extra:202:4}{}
		B. Nachtergaele and R. Sims, \emph{Lieb--Robinson bounds in quantum
			many-body physics}, in R. Sims and D. Ueltschi (eds.), \emph{Entropy
			and the Quantum}, Contemporary Mathematics 529, American Mathematical
		Society (2010), 141--176, Section 3.1, Theorem 3.1 and its proof
		(norm convergence of bounded-interaction dynamics on local observables).
		\url{https://arxiv.org/pdf/1004.2086v1}.
		{\small\textit{Locality input for the thermodynamic-limit argument;
				the Gaussian-smearing passage of the KMS identity is given in the text.
				Consulted 24 September 2026.}}
		\item[\textnormal{[E5]}]\hypertarget{extra:202:5}{}
		D. Heson, S. Starr, and J. Thornton, \emph{Violation of ferromagnetic
			ordering of energy levels in spin rings for the singlet},
		arXiv:2307.12773v5, 7 January 2025, abstract and Section 2.
		\url{https://arxiv.org/html/2307.12773v5}.
		{\small\textit{The manuscript distinguishes its largely numerical
				ordering claim from a rigorous singlet-sector result; neither is
				used as a proof of the lattice conjecture here. Consulted
				24 September 2026.}}
		% END RESEARCH REFERENCES 202
	\end{itemize}
\end{document}
